\uuid{GpMF}
\exo7id{3869}
\titre{exo7 3869}
\auteur{quercia}
\organisation{exo7}
\datecreate{2010-03-11}
\isIndication{false}
\isCorrection{true}
\chapitre{Continuité, limite et étude de fonctions réelles}
\sousChapitre{Continuité : théorie}
\cpge{mpsi-mp2i.fonctions-valeurs-reelles-complexes.a, pcsi.techniques-calcul-fonctions.a, ptsi.techniques-calcul-fonctions.a, mpsi-mp2i.fonctions-valeurs-reelles-complexes.c, pcsi.techniques-calcul-fonctions.c, ptsi.techniques-calcul-fonctions.c, bcpst1.limites-continuite.c, bcpst1.fonctions-reelles-usuelles.a}
\module{Analyse}
\niveau{L1}
\difficulte{}

\contenu{
\question{
Trouver toutes les fonctions $f : {\R} \to {\R}$ continues telles
que : $\forall\ x,y \in \R,\ f(x+y)f(x-y) = f^2(x)f^2(y)$.
}
\reponse{
Si $f$ n'est pas identiquement nulle, alors $f(0) = \pm1$ et $f$ est paire,
de signe constant.

Par récurrence, $\forall\ p \in \N,\ f(px) = \pm f^{p^2}(x)$
$ \Rightarrow $ par densité, $f(x) = \pm\lambda^{x^2}$.
}
}
